\chapter{AI-Informed Solvation Engineering of Liquid Thermogalvanic
Electrolytes}

\section{Introduction}

Thermogalvanic cells convert a temperature difference into a continuous
electrochemical potential through the temperature dependence of a redox
reaction. For a simple redox couple in a liquid electrolyte,
\begin{equation}
S_e=\frac{\Delta E}{\Delta T}
=\frac{\Delta S_{\mathrm{rc}}}{nF}
=\frac{1}{nF}
\left(
\frac{\mathrm{d}G_{\mathrm{solv}}^{\mathrm{O}}}{\mathrm{d}T}
-
\frac{\mathrm{d}G_{\mathrm{solv}}^{\mathrm{R}}}{\mathrm{d}T}
\right),
\end{equation}
where $G_{\mathrm{solv}}^{\mathrm{O}}$ and
$G_{\mathrm{solv}}^{\mathrm{R}}$ are the solvation free energies of the
oxidized and reduced species. A solvent that perturbs the two redox states
unequally can therefore enlarge the solvation entropy difference without
changing the electrode architecture \citep{chen2023solvation}.

The apparent simplicity of this strategy hides two coupled difficulties.
First, thermopower is not controlled by a single solvent property. It emerges
from differences in the interaction energies, coordination geometries, packing,
orientation, and collective dynamics of solvent molecules around two redox
ions. Second, the experimental database is small. The scarcity of $S_e$
measurements for thermogalvanic electrolytes prevents the construction of a
large and consistent training set, and the difficulty increases when
ion--ion and ion--solvent interactions are included in the prediction. A
direct regression of $S_e$ from molecular properties is therefore readily
overfitted. Data-driven work on ionic thermoelectrics and thermal-energy
materials has reported the same dependence on dataset size and physically
relevant features \citep{wu2024unlocking,li2022mlthermal}.

Here, DFT calculations were combined with machine learning to screen solvent
molecules for the $[\mathrm{Fe(CN)}_6]^{4-/3-}$ electrolyte. Ion--solvent
interaction features were first parameterized into a descriptor correlated
with $S_e$. The descriptor was then used for high-throughput solvent screening,
followed by experimental evaluation and refinement of the model.

\section{Correlation between Ion--Solvent Interactions and Thermopower}

The benchmark system was the aqueous
$[\mathrm{Fe(CN)}_6]^{4-/3-}$ redox couple. A useful co-solvent must compete
with water for access to the redox ions while creating different local
environments around the two charge states. Density functional theory (DFT)
calculations were used to quantify ion--solvent interaction energies, bond
lengths, and bond angles. These quantities represent interaction strength, the
dimensions of the local coordination environment, and the packing or
orientation of solvent molecules around the ions, respectively.

For ten solvents with reported thermopower values, the differences in these
interaction features between $[\mathrm{Fe(CN)}_6]^{4-}$ and
$[\mathrm{Fe(CN)}_6]^{3-}$ were correlated with $S_e$. No single feature was
sufficient, which is consistent with thermopower being a collective solvation
property. Partial least-squares dimensionality reduction was consequently used
to combine the interaction features into a single descriptor, $\alpha$. This
descriptor correlated strongly with thermopower and varied monotonically with
$S_e$, providing an intermediate target that was more data-efficient and more
physically interpretable than a direct solvent-to-thermopower mapping
.

The parameterization also revealed an asymmetry between the two redox states.
Features associated with the more highly charged
$[\mathrm{Fe(CN)}_6]^{4-}$ ion showed stronger correlations with
thermopower. Its stronger ion--dipole interactions produce a more ordered
solvation shell and make its entropy more sensitive to temperature-induced
reorganization. Promising solvents should therefore not simply interact
strongly with both ions; they should enlarge the difference between their
local responses.

\section{High-Throughput Solvent Screening}

High-throughput calculations generated $\alpha$ values for 64 candidate
solvents. Each molecule was represented by ten molecular and physicochemical
features: molecular weight, polar surface area, molecular complexity,
effective three-dimensional rotor count, density, melting point, viscosity,
dielectric constant, dipole moment, and donor number. Dielectric constant,
dipole moment, and donor number were included explicitly because of their
direct relevance to ion solvation.

The limited dataset made precise numerical regression unstable. The problem
was therefore reformulated as interval prediction using a discretized-target
gradient boosting classifier (DT-GBC). A Gaussian mixture model first divided
the calculated $\alpha$ values into low, medium, and high classes; a gradient
boosting classifier then predicted the class of an unseen solvent. Converting
an exact-value problem into a chemically useful ranking problem reduced the
effective learning difficulty. The model classified 10 of 13 test samples
correctly (76.9\% accuracy). Repeated ten-fold cross-validation produced a mean
accuracy of 0.64 under changes in the small and imbalanced data splits,
quantifying uncertainty that would be hidden by a single train--test partition
.

SHAP analysis identified dipole moment as the most influential feature,
followed by donor number, polar surface area, and dielectric constant. These
features describe complementary aspects of polarity: charge separation,
electron donation to an ion, accessibility of polar groups, and dielectric
response to charge. Their collective importance established polarity-related
features as the most important variables in the model. This result does not
imply that thermopower varies monotonically with any one polarity descriptor.

\section{Experimental Validation and Model Refinement}

The initial predictions were tested using six solvents. Triethyl phosphate
(TEP), N,N-dimethylformamide (DMF), and
ethylenediaminetetraacetic acid followed the predicted trend, whereas glycine,
glucose, and erythritol did not. These failures exposed the boundary of the
initial representation: it described simple two-body ion--solvent interactions
but did not fully capture steric shielding, molecular symmetry, dissociation,
or extended hydrogen-bond networks.

The discrepancies were used to add a four-dimensional binary filter based on
carbon-chain length, steric hindrance, dissociation energy, and hydrogen-bond
donor/acceptor count. The combined procedure correctly classified 15 of 16
evaluated solvents, corresponding to 93.8\% accuracy. Thus, the model produced
a testable ranking, failed predictions exposed missing physics, and targeted
experiments improved the representation.

The refined framework identified TEP, tetrahydrofuran, dimethylacetamide,
methylsulfonylmethane, and 3-methyl-2-oxazolidinone as effective candidates,
with thermopower enhancements of approximately 42\% under the screening
conditions. Molecular dynamics simulations further distinguished effective
and unsuccessful solvents. Effective molecules such as DMF and TEP produced
clearly different solvation structures around the two redox ions, whereas
glucose caused only a small difference. The measured thermopower was therefore
consistent with asymmetric solvation reorganization rather than solvent
polarity alone.

\section{Solvent Modification of a Thermosensitive Electrolyte}

The identified solvents were subsequently evaluated in a guanidinium chloride
(GdmCl)-based electrolyte. A 10 wt\% GdmCl solution produced
$S_e=2.52$ mV K$^{-1}$. Adding TEP or DMF increased the thermopower to
3.35 and 3.26 mV K$^{-1}$, respectively. Composition optimization produced a
maximum of 4.0 mV K$^{-1}$ at 20 wt\% TEP and 20 wt\% GdmCl.
UV--visible spectra showed that TEP affected
$[\mathrm{Fe(CN)}_6]^{4-}$ more strongly than
$[\mathrm{Fe(CN)}_6]^{3-}$, consistent with the larger sensitivity of the
high-valence ion found in the DFT and molecular-dynamics results.

